% Chapter source (edited directly; this is the master copy of the text)
\chapter{CONCLUSION AND FUTURE WORK}

\section{Answers to the research questions}

\textbf{RQ1, feasibility at the edge: yes.} On a Raspberry Pi 5 with a Hailo-8, the deployment setting (n model with tiles) ran 15 minutes in real time:

\begin{itemize}
\tightlist
\item
  30 FPS with 47 ms latency
\item
  4.1 W measured on the board, plus about 2.5 W for the accelerator
\item
  peak SoC temperature 60 °C, no throttling
\end{itemize}

The pipeline is limited by the accelerator, not the CPU. The m model with tiles reaches only 7.9 FPS, enough for analysis but not for a live camera without dropping frames.

\textbf{RQ2, coverage: substantially improvable without retraining.} Tiling the back of the room raised the share of SCB students found from 55\% to 72\%, and in the far third from 25\% to 66\%. It more than doubled the students found in a dense study hall. Removing same-skeleton duplicates and anchoring seats at the head brought the inferred class size close to the real head count on the lecture (26 for 28 students), but it still overcounts the exam hall (29 for 23 students and a teacher) and the dense study hall (59 for about 50).

In a small room filmed at close range (clip\_5), seats multiplied too, to 13-16 for 10 students, so the class-size estimate is a rough guide, not a head count.

\textbf{RQ3, the per-student decision: calibration beats benchmark accuracy.} Classifiers trained on SCB rank SCB actions well (within-image AUC 0.83) but transfer badly:

\begin{itemize}
\tightlist
\item
  \textbf{Lecture clip:} they call 50-82\% of an attentive lecture distracted.
\item
  \textbf{Human labels:} they wrongly flag 84-159 of 262 focused moments.
\item
  \textbf{Cause:} SCB has no ``listening'' label, so they learn that a raised head means distracted.
\item
  \textbf{A second dataset:} adding ClaBehavior, which has that head-up class, lowers the lecture share only to 44-51\% for trees and MLP, and a model trained on one dataset reaches only AUC 0.70-0.81 on the other.
\end{itemize}

The calibrated rules have very few false alarms (specificity 95\% on human labels in the deployment setting; 1-2\% on the lecture). They first missed most real distraction (11\% on the staged session), because the ear-based head yaw collapses in profile. A second, body-relative head-turn measure raised that to 56\% with both models, with no new false alarms there and a cost of 7-10 points on the exam hall. On a 10-student class labelled after all thresholds were fixed (clip\_5), the rules on the deployment setting kept 89\% specificity. Phone use remains invisible to pose.

\textbf{RQ4, the class indicator: yes, with honest width.} Seat memory, zone weighting and a Wilson interval around the shown value give a share that always lies inside its range. Without the real head count the range is wide (15-44 points on the clips), which reflects how little a sample that small can say. Giving \texttt{-\/-class-size} narrows it.

\section{Future work}

\begin{enumerate}
\def\labelenumi{\arabic{enumi}.}
\tightlist
\item
  \textbf{A held-out test at classroom scale.} Run the 9-person scripted session (research/staged/script\_v2.csv). clip\_5 confirmed the rules' specificity on a class nothing was tuned on, but its recall still rests on one two-person session.
\item
  \textbf{Classroom-specific training labels.} ClaBehavior showed that a public dataset with a head-up class is not enough. The labelling and staged-session tools can grow a dataset from the target classrooms, with attentive, head-up students labelled over time. A learned decision trained on it, with per-seat baselines as features, could combine the classifiers' sensitivity with the rules' specificity.
\item
  \textbf{Phone use.} Detecting phones needs an object detector. On the Hailo-8 it could run on the same frames, at a cost in frame rate. Only a per-seat ``phone visible'' flag would leave the device.
\item
  \textbf{Longer field deployments.} Weeks of use would show drift in calibration, the effect of lighting and seating changes, and whether teachers find the indicator useful. That needs a study with consent, designed with the school.
\item
  \textbf{Live m with tiles.} A leaky queue or a lower camera rate would let the most accurate setting run live without growing latency.
\item
  \textbf{Keypoint-based geometry.} Anchor seats on the shoulders and measure torso rotation against the visible upper body instead of the box. That removes the two failures seen at close range in clip\_5: seats multiplying and a full-body box reading as a turned torso.
\item
  \textbf{Teacher-facing evaluation.} The real test of a class-level indicator is whether it helps teaching, which needs a user study rather than more accuracy metrics.
\end{enumerate}
